\documentclass[12pt, oneside, a4paper]{ctexart}
\usepackage{amsmath}
\usepackage{amssymb}
\usepackage{caption}
\usepackage{ctex}
\usepackage{enumerate}
\usepackage{float}
\usepackage{tocloft}
\usepackage[{top=2.54cm, bottom=2.54cm, left=3.18cm, right=3.18cm}]{geometry}
\usepackage{graphicx}
\usepackage[breaklinks, colorlinks=false, hidelinks]{hyperref}
\usepackage{mathptmx}
\usepackage{xcolor}
\usepackage{xeCJK}

\ctexset{
    section={   
        name={,、},
        number={\heiti\chinese{section}},
        format=\heiti\raggedright\zihao{3}\bfseries,
        aftername={\hspace{0.5em}},
        beforeskip=1ex,
        afterskip=1ex,
    },
    subsection={
        name={,.},
        number={\arabic{section}.\arabic{subsection}},
        format=\songti\raggedright\zihao{4}\bfseries,
        aftername={\hspace{0.5em}},
        beforeskip=0ex,
        afterskip=0ex,
    },
    subsubsection={
        name={,.},
        number={\arabic{section}.\arabic{subsection}.\arabic{subsubsection}},
        format=\songti\raggedright\zihao{-4}\bfseries,
        aftername={\hspace{0.5em}},
        beforeskip=0ex,
        afterskip=0ex,
    },
}
\setlength{\cftbeforesecskip}{1ex plus 0.1ex}
\setlength{\cftbeforesubsecskip}{0ex plus 0.05ex}
\setlength{\cftbeforesubsubsecskip}{0ex plus 0.05ex}
\pagestyle{plain}
\renewcommand{\figurename}{\small{图}}
\setmainfont{Times New Roman}
\setCJKmainfont[AutoFakeBold={2}]{SimSun}
\setCJKfamilyfont{zh}{SimSun}
\setCJKfamilyfont{ja}{BIZ UDMincho}
\xeCJKDeclareCharClass{CJK}{
  "24EA,
  "2460->"2473,
  "3251->"32BF,
  "24FF,
  "2776->"277F,
  "24EB->"24F4
}



%----------分割线----------%



\begin{document}



\thispagestyle{empty}
\begin{center}{
    \fontsize{26pt}{26pt}\selectfont
    \vspace*{2.75\baselineskip}
    \includegraphics[width=0.6\textwidth]{pics/tongji-university-logo-1024px.png}\\
    [\baselineskip]
    \heiti{轨道车辆设计作业（二）}\\
    \fontsize{18pt}{24pt}\selectfont
    \vspace*{2\baselineskip}
    \songti{\begin{tabular}{l}
            {\heiti{}完成日期：}2026年5月29日 \\
            {\heiti{}报告作者：}肖哲夫~~2352670
        \end{tabular}
        \vspace{0.25cm}
        \\\textcolor{gray}{\textit{Made with }\LaTeX}}
    }
\end{center}
\newpage{}
\setcounter{page}{1}
\begin{center}
    \heiti\zihao{2} 轻量化设计的系统思维
\end{center}
\vspace{-1cm}
\begin{center}
    \heiti\zihao{2} ——材料 - 结构的全链条优化
\end{center}
\vspace{1ex}
\section{\heiti{概述}}
\par{}
轻量化设计是现代工程领域中最具跨学科特征的技术方向之一\textsuperscript{\textcolor{red}{\cite{ref1}}}，阅读伯恩德·克莱恩教授所著的《轻量化设计——计算基础与构件结构》，我系统了解了轻量化设计的理论基础、深刻理解了轻量化作为一门工程科学的整体思维框架。
\par{}
\quad{}
\par{}
\section{\heiti{核心思想与主要内容}}
\par{}
本书从材料轻量化、形状轻量化、制造轻量化等几个维度，从平衡技术性能、经济成本与可持续发展并寻找其最优解的优化轻量化理念出发，强调使用轻量化的实际收益，如降低能耗、提升性能等来衡量轻量化措施的价值。
\par{}
从全书内容来看，这26章的编排存在鲜明的逻辑递进关系，前7章从轻量化的目标、问题结构、方法与工具、轻量化构造、材料选择准则、轻量化材料、轻量化设计原则等方面，详细阐述了作者的设计理念，尤其是材料设计准则一章，将看似简单的问题从量化视角出发，如用轻量化系数LBK来衡量轻量化程度。
\par{}
第8至16章从弹性理论等出发，针对薄壁型杆、开口型杆、桁梁型材、箱型材等多种形状的结构逐步构建分析轻量化结构的理论工具。尤其是力流法的引入，有助于设计者直观理解力在结构中的传播途径，从而针对性优化材料配置。
\par{}
第17至22章则聚焦特殊构造与连接技术，如三明治单元、加固性设计等。其中板材的凸起分析一章占用了极大篇幅，也体现了稳定性问题在薄壁轻量化结构中的重要地位。
\par{}
最后4章则以结构优化、疲劳寿命、可靠性和结构声学四个主题代表了现代轻量化设计的前沿方向。
\par{}
\quad{}
\par{}
\section{\heiti{感想}}
\par{}
纵观全书，我认为该书有三个鲜明的特点：
\par{}
\symbol{"2460}理论与实践高度统一。本书的最后近一百页全是练习，均通过案例来加深读者的理解，并配有详细的解题过程，而不只是一条算式作为答案，对于初学者较为友好。虽然时间有限，我并未深入探究这些习题，但我认为这些习题非常值得慢慢品鉴。
\par{}
\symbol{"2461}成本意识贯彻始终。从头几章，尤其是第2章提出的成本模型开始，作者始终秉持着这样的理念：工程的本质是在约束条件下求解最优方案，忽视成本的轻量化不是真正的工程。结合轨道车辆设计与工程项目管理等课程所学内容，我意识到这些理念的确已被贯彻到大多数设计的方方面面，将来做设计时我也应注意成本方面。
\par{}
\symbol{"2462}本书覆盖了从材料微观性能（如弹性模量、断裂韧性）到宏观结构（杆、梁、壳体、三明治板），再到结构系统性能（疲劳寿命、可靠性、声学行为）的完整知识链条，培养读者从整体角度理解轻量化问题的能力。
\par{}
对于车辆工程专业而言，轻量化设计的价值尤为突出。不论是汽车还是轨道车辆，不论是传统燃油车辆、新能源汽车还是电力牵引轨道车辆，对节约能源与制造成本均有一定要求，同时对强度、寿命等指标也存在要求。因此该书不仅提供了设计轻量化结构的一般方法，还培养了读者面向约束条件求解最优方案的工程思维，有利于我们面向生产实践进行轻量化结构的设计研究。
\par{}
\quad{}
\par{}
\begin{flushleft}
    \begin{small}
        \begin{thebibliography}{99}
            \bibitem{ref1}Bernd Klein. 轻量化设计  计算基础与构件结构  第2版[M]. 北京: 机械工业出版社, 2016: 2.
            \vspace{-3mm}
        \end{thebibliography}
    \end{small}
\end{flushleft}



\end{document}